% Queue full-width controlled figures before their discussion so they remain
% near Section IV rather than drifting through the Adam section.

\begin{figure*}[!t]
  \centering
  \includegraphics[width=\textwidth]{figures/analysis024/gd_loss_projected_change_spectral_clean.pdf}
  \caption{Training loss and curvature diagnostics under full-batch GD.
  Columns identify dataset/model; colors identify learning rate.
  Top: unsmoothed full-objective loss on logarithmic axes.
  Bottom: $\eta\widehat\kappa_t$ (solid, complete 50-pair trailing means)
  and retained $\eta\lambda_{\max}$ measurements (dashed with circles);
  the dotted line marks two. Spectra use the complete MLP objective,
  5,000/50,000 ResNet images, and 2,500/59,680 Transformer sequences.
  Loss uses parameter state $s$; a spectrum at state $s$ is placed at
  gradient index $t=s+2$, the starting state for $\widehat\kappa_t$.
  The final spectrum has no recorded gradient pair.}
  \label{fig:gd-loss}
\end{figure*}

\begin{figure*}[!t]
  \centering
  \includegraphics[width=\textwidth]{figures/analysis024/gd_reversal_norm_retention.pdf}
  \caption{Constructive/destructive alignment and cumulative motion for
  the conditions in Fig.~\ref{fig:gd-loss}. Rows show consecutive cosine,
  gradient norm (logarithmic scale), and normalized displacement.
  Cosines and norms use complete 50-observation trailing means;
  $R_t$ is the unsmoothed ratio over the prefix through update $t$.
  Colors and column ordering match Fig.~\ref{fig:gd-loss}.}
  \label{fig:gd-alignment-displacement}
\end{figure*}

\begin{figure*}[!t]
  \centering
  \includegraphics[width=\textwidth]{figures/analysis035/gd_backward_compact.pdf}
  \caption{Backward gradient alignment at the common endpoint
  $\theta_{2256}$, corresponding to $t=2257$.
  Each point is one exact $\cos(g_t,g_{t-k})$ with fixed endpoint.
  Top: lags 1--256. Bottom: the same values at lags 1--32.
  Rates are $0.04$ (MLP), $0.1$ (ResNet) and $0.2$ (Transformer).
  These separately recomputed 256-update GD continuations start from
  retained state 2000, with one initialization, no smoothing and no
  endpoint averaging.}
  \label{fig:gd-backward}
\end{figure*}

\begin{figure*}[!t]
  \centering
  \includegraphics[width=\textwidth]{figures/analysis037/sgd_batch_size_compact.pdf}
  \caption{Mean sampled-gradient cosine over 199 adjacent pairs during
  200-update continuations. Blue circles start at initialization;
  orange squares start after 2,000 GD updates. Each line averages
  repetitions within initialization, then weights the three
  initialization blocks equally; shading is their range.
  Hollow rightmost markers denote full GD. ResNet batches contain
  whole fixed 1,000-image normalization groups; Transformer batches
  count prepared sequences. Batch axes are logarithmic.}
  \label{fig:sgd-moving-geometry}
\end{figure*}
